% Part: sets-functions-relations
% Chapter: infinite
% Section: dedekinds-proof
%
\documentclass[../../../include/open-logic-section]{subfiles}

\begin{document}

\olfileid[fa-IR]{sfr}{infinite}{dedekindsproof}

\olsection[«برهان» ددکیند]{«برهان» ددکیند بر وجود یک
مجموعهٔ نامتناهی}

در این فصل، اعداد طبیعی را برحسب جبرهای ددکیند به شیوه‌ای مجموعه‌شناختی
بررسی کردیم. در \olref[arith][ref]{sec} دربارهٔ اهمیت فلسفیِ حسابی‌سازیِ
تحلیل (در کنار موضوعات دیگر) تأمل کردیم. اکنون باید دربارهٔ اهمیت آنچه
اینجا به دست آورده‌ایم تأمل کنیم.

در سراسر \olref[sfr][arith][]{chap} اعداد طبیعی را مفروض گرفتیم و با
استفاده از آن‌ها اعداد صحیح، گویا و حقیقی را به‌صراحت ساختیم. در این فصل
ساختی صریح از اعداد طبیعی عرضه نکرده‌ایم. فقط نشان داده‌ایم که
\emph{هر مجموعهٔ نامتناهیِ ددکیندی که داده شود}، می‌توانیم مجموعه‌ای تعریف
کنیم که دقیقاً چنان رفتار کند که از~$\Nat$ می‌خواهیم.

پس آشکار است که نمی‌توانیم مدعی شویم به پرسشی مابعدالطبیعی، مانند اینکه
\emph{کدام اشیا همان اعداد طبیعی‌اند}، پاسخ داده‌ایم. اما این نکتهٔ خوبی
است. گذشته از این، در \olref[sfr][arith][ref]{sec} تأکید کردیم که اشتباه
است تعریف $\Real$ به‌صورت مجموعهٔ برش‌های ددکیند را به‌جای یک قرارداد
سودمند، \emph{کشف} بدانیم. مشاهدهٔ اساسی این است که برش‌های ددکیند خواص
ریاضیِ اصلیِ اعداد حقیقی را متجلی می‌کنند. در اینجا نیز چنین است: مشاهدهٔ
اساسی این است که \emph{هر} جبر ددکیند خواص ریاضیِ اصلیِ اعداد طبیعی را
متجلی می‌کند. (در واقع، ددکیند با اثبات اینکه همهٔ جبرهای ددکیند
\emph{هم‌ریخت} هستند این نکته را قاطعانه نشان داد
(\citeyear[قضیه‌های 132--3]{Dedekind1888}). پس جای شگفتی نیست که بسیاری
از «ساختارگرایان» معاصر ددکیند را پیشگام خود می‌دانند.)

افزون بر این، نشان داده‌ایم چگونه نظریهٔ اعداد طبیعی را در یک نظریهٔ
طبیعی و سادهٔ مجموعه‌ها بنشانیم؛ نظریه‌ای که خود هنوز نسبتاً غیرصوری
است، اما (ظاهراً) اعداد طبیعی را مفروض نمی‌گیرد. بنابراین شاید در مسیر
تحقق طرح بلندپروازانهٔ خودِ ددکیند باشیم؛ طرحی که او چنین توضیح داد:
\begin{quote}
	در علم، هیچ امرِ اثبات‌پذیری را نباید بی‌برهان پذیرفت. هرچند این مطالبه
	معقول می‌نماید، نمی‌توانم آن را حتی در جدیدترین شیوه‌های بنیان‌گذاریِ
	ساده‌ترین علم برآورده‌شده بدانم؛ یعنی آن بخش از منطق که به نظریهٔ اعداد
	می‌پردازد. وقتی از حساب (جبر، تحلیل) صرفاً همچون بخشی از منطق سخن
	می‌گویم، مقصودم این است که مفهوم عدد را یکسره مستقل از مفاهیم یا شهودهای
	مکان و زمان می‌دانم---بلکه آن را محصول بی‌واسطهٔ قوانین ناب اندیشه
	می‌دانم.
	\citep[پیشگفتار]{Dedekind1888}
\end{quote}
اندیشهٔ جسورانهٔ ددکیند چنین است. همین اکنون نشان داده‌ایم چگونه اعداد
طبیعی را تنها با استفاده از نظریهٔ طبیعی مجموعه‌ها بسازیم. در
\olref[sfr][arith][]{chap} دیدیم که با داشتن اعداد طبیعی و مقداری نظریهٔ
مجموعه‌ها چگونه اعداد حقیقی را بسازیم. پس شاید «حساب (جبر، تحلیل)» «صرفاً
بخشی از منطق» از آب درآید (با معنای گسترده‌ای که ددکیند برای واژهٔ «منطق»
در نظر داشت).

اندیشه همین است. اما لحظه‌ای درنگ کنید. ساخت ما از یک جبر ددکیند (بدیل ما
برای اعداد طبیعی) مشروط به وجود یک مجموعهٔ نامتناهیِ ددکیندی است. (کافی
است به \olref[sfr][infinite][dedekind]{thm:DedekindInfiniteAlgebra}
بازگردید.) مگر آنکه بتوان وجود یک مجموعهٔ نامتناهیِ ددکیندی را از راه
«منطق» یا «قوانین ناب اندیشه» اثبات کرد، طرح متوقف می‌شود.

پس آیا \emph{می‌توان} وجود یک مجموعهٔ نامتناهیِ ددکیندی را با «قوانین ناب
اندیشه» اثبات کرد؟ تلاش ددکیند چنین بود:
\begin{quote}
  قلمرو افکار خود من، یعنی تمامیتِ $S$، متشکل از همهٔ چیزهایی که می‌توانند موضوع
  اندیشهٔ من باشند، نامتناهی است. زیرا اگر $s$ نشانهٔ عضوی از~$S$ باشد،
  آنگاه اندیشهٔ $s'$ مبنی بر اینکه~$s$ می‌تواند موضوع اندیشهٔ من باشد،
  خود عضوی از~$S$ است. اگر این را تصویری به‌صورت $\phi(s)$ از عضو~$s$ بدانیم،
  آنگاه \dots~$S$ [به معنای ددکیندی] نامتناهی است، و این همان چیزی بود
  که باید ثابت می‌شد.
	\citep[\S66]{Dedekind1888}
\end{quote}
یافتن چنین مطلبی در میانهٔ کتابی که بخش عمده‌اش از برهان‌های ریاضیِ بسیار
دقیق تشکیل شده، کاملاً شگفت‌آور است. دو نکته شایستهٔ ذکرند.

نخست: این «برهان» به‌سختی خصلتی دارد که امروزه آن را «ریاضی» بشناسیم. این
برهان از اشیای روان‌شناختی (اندیشه‌ها) سخن می‌گوید، آن هم اشیایی که صرفاً
\emph{ممکن} هستند.

دوم: دست‌کم با نحوه‌ای که ما جبرهای ددکیند را عرضه کرده‌ایم، این «برهان»
کاستیِ فنیِ سرراستی دارد. اگر استدلال ددکیند موفق باشد، تنها ثابت می‌کند
که بی‌نهایت چیز وجود دارد (به‌طور مشخص، بی‌نهایت اندیشه). اما ددکیند همچنین
باید دلیلی به ما بدهد تا~$S$ را یک \emph{مجموعهٔ} واحد، که بی‌نهایت
!!{element}s دارد، بدانیم، به‌جای آنکه~$S$ را \emph{چیزهایی} (به صیغهٔ
جمع) تلقی کنیم.

این واقعیت که ددکیند در اینجا شکافی ندید شاید نشان دهد کاربرد واژهٔ
«تمامیت» نزد او دقیقاً با کاربرد \emph{ما} از واژهٔ «مجموعه» متناظر نیست.
\footnote{در واقع، دلایل دیگری نیز داریم که چنین نبوده است؛ بنگرید به
\citet[ص.~23]{Potter2004}.} اما این امر چندان شگفت‌آور نیست. طرحی که در دو
فصل گذشته دنبال کردیم---«ساختن» اعداد طبیعی و از آن‌ها «ساختن» اعداد صحیح،
حقیقی و گویا---یکسره به شیوه‌ای ساده‌انگارانه انجام شده است. هر زمان که به
این مجموعه یا آن مجموعه نیاز داشته‌ایم، آن را مفروض گرفته‌ایم، بی‌آنکه
اصول کلیِ بسیاری دربارهٔ اینکه دقیقاً کدام مجموعه‌ها و چه هنگام وجود دارند
وضع کنیم. اما می‌دانیم به \emph{برخی} اصول کلی نیاز داریم؛ وگرنه به
پارادوکس راسل گرفتار خواهیم شد.

زمان آن رسیده است که از ساده‌انگاریِ خود فراتر برویم.

\end{document}
